\documentclass[10pt,a4paper]{amsart}
\usepackage[margin=19mm]{geometry}
\usepackage{amsmath,amssymb,mathtools}
\usepackage[scr=rsfs,cal=euler]{mathalfa}
\usepackage{xfrac}
\usepackage[unicode,hidelinks,bookmarks=false]{hyperref}
\newcommand{\ie}{i.e.\ }
\newcommand{\cf}{cf.\ }
\newcommand{\resp}{respectively\ }
\newcommand{\Subsection}{\subsection}
\providecommand{\nobreakdash}{\nobreak\hbox{-}}
\setlength{\emergencystretch}{2em}
\multlinegap0pt
\usepackage{bm}
\usepackage{bbm}
\usepackage{dsfont}
\usepackage{xspace}
\usepackage{cancel}
\usepackage{mathtools}
\usepackage{amsthm}
\makeatletter
\@ifundefined{theorem}{\newtheorem{theorem}{Theorem}}{}
\@ifundefined{lemma}{\newtheorem{lemma}[theorem]{Lemma}}{}
\makeatother
\makeatletter
\def\E{\mathbb{E}}
\def \R{\mathbb {R}}
\expandafter\ifx\csname remerciements\endcsname\relax
\newenvironment*{remerciements}{
\renewcommand*{\abstractname}{Remerciements}
\begin{abstract}
}{\end{abstract}}
\fi
\renewcommand{\sup}{sup\,}
\def\vp{\varepsilon}
\makeatother
\title[Scaling limit of first passage percolation geodesics on plan]{Scaling limit of first passage percolation geodesics on planar maps}
\author{Emmanuel Kammerer}
\date{Source: arXiv:2412.02666v3 (CC BY 4.0)}
\begin{document}
\maketitle

\noindent\emph{Source and licence.} Scaling limit of first passage percolation geodesics on planar maps. Version arXiv:2412.02666v3 (CC BY 4.0), distributed under the Creative Commons Attribution 4.0 International licence, as recorded in the arXiv licence metadata for this exact version and shown on the abstract page https://arxiv.org/abs/2412.02666v3. Full rights notice: licenses/arxiv-2412.02666-CC-BY-4.0.txt.\par\smallskip

\noindent\textit{Editorial scope.} Counted unit: the Lemma whose source label is \texttt{lemme couplage} in the article (source0.tex lines 2049--2062, source label \texttt{lemme couplage}), delivered together with the article's own complete original proof of it (source0.tex lines 2064--2142, one \texttt{proof} environment of 79 lines). No mathematics is written, repaired or re-derived: statement and proof are the article's own text line for line. Required context retained uncounted: eq majoration g discret (lines 1823--1825); lemme martingale (lines 1836--1852); lemme petits sauts de P (lines 1855--1866). Every definition, display and label of those lines is retained unchanged. Omitted from this package and not redistributed: 1--1822, 1826--1835, 1853--1854, 1867--2048, 2063--2063, 2143--2845. Nothing inside a retained excerpt is omitted or altered: every retained body is delivered exactly as the article's lines are written, so no omission receipt of any kind accompanies this package. Citations: the retained excerpts carry no citation command at all, so no bibliography entry of the article is transcribed and no reference is redistributed. Reference adaptation: none was needed and none was made; every reference inside the delivered unit resolves inside it, so no unresolved reference or citation is produced. Printed slips kept literal: no printed word, symbol, hypothesis or number of the counted statement or of its proof was altered. Layout and class: the article's own class is \texttt{article} with packages including babel, inputenc, fontenc, amsmath, amssymb, fullpage, xy, tikz, enumerate, mathtools, graphicx, color, hyperref, pgf; the wrapper is the lane's pinned \texttt{amsart} editorial wrapper at 10pt on A4, which restates the theorem-like environments the retained text needs and adds the article's own macro definitions that the retained text needs (\texttt{\textbackslash E}, \texttt{\textbackslash R}, \texttt{\textbackslash sup}, \texttt{\textbackslash vp}), with extra wrapper packages (bm, bbm, dsfont, xspace, cancel, mathtools). The delivered numbering is the unit's own. Assets: no retained span contains a figure environment, a graphics-inclusion command, a PDF-page inclusion, a table or a TikZ picture; no image file is copied into this package, the package declares no asset dependency and no image file is redistributed with it. Licence: version arXiv:2412.02666v3 of this article is distributed under the Creative Commons Attribution 4.0 International licence, as recorded in the arXiv licence metadata for this exact version and shown on the abstract page https://arxiv.org/abs/2412.02666v3. A full rights notice is delivered at licenses/arxiv-2412.02666-CC-BY-4.0.txt. Characters: every retained excerpt is pure ASCII, so the delivered engine, XeLaTeX with its default Latin Modern text font, reports no missing character.
\begin{equation}\label{eq majoration g discret}
	g_t^{(\ell),T}(x,z,u) \le 3 \vert z-1 \vert.
\end{equation}

\begin{lemma}\label{lemme martingale}
	Let $\mathcal{F}^{(\ell),T}$ be the filtration defined by
	$$
	\forall t \in [0,T], \qquad
	\mathcal{F}^{(\ell),T}_t
	= \sigma\left(
	\widetilde{P}((T-s)-), U_{T-s}^{(\ell)}; \ s \in [0,t]
	\right).
	$$
	Then for all $\vp>0$, for all $v\in \R$, the process $(M^{(\ell),T,\vp}_t)_{t \in [0,T]}$ defined by
	$$
	\forall t \in [0,T],\qquad
	M^{(\ell),T,\vp}_t
	=\int_0^t\int_{1-\vp}^{1+\vp} \int_0^1 g_s^{(\ell),T} ( X^{(\ell),T}_{s-}(v),z,u) N^{(\ell),T}(ds,dz,du)
	$$
	is a martingale with respect to the filtration $\mathcal{F}^{(\ell),T}$.
\end{lemma}

\begin{lemma}\label{lemme petits sauts de P}
	We have the uniform convergence
	$$
	\sup_{\ell \ge 1}
	\E^{(\ell)}
	\left(
	\sum_{s \in[0,T]} {\bf 1}_{\widetilde{P}(s)/\widetilde{P}(s-) \in [1-\vp,1+\vp]\setminus \{1\}}
	\left(\frac{\Delta \widetilde{P}(s)}{\widetilde{P}(s-)}\right)^2 
	\right)
	\mathop{\longrightarrow}\limits_{\vp \to 0}0.
	$$
\end{lemma}

\begin{lemma}\label{lemme couplage}
For every sequence of real numbers $(v^{(\ell)})_{\ell \ge 1}$, there exist couplings of $(X^{(\ell),T}_t(v^{(\ell)}))_{t \in [0,T]}$ with $(X^{(\ell),T,\vp}_t(v^{(\ell)}))_{t \in [0,T]}$ for all $\ell \ge 1$ and for all $\vp >0$, denoted by $X^{(\ell),T}(v^{(\ell)}),\overline{X}^{(\ell),T,\vp}(v^{(\ell)})$ such that
$$
\sup_{\ell \ge 1} \E^{(\ell)} 
\left(\sup_{t\in [0,T]} 
\left|
X^{(\ell),T}_t(v^{(\ell)})
-
\overline{X}^{(\ell),T,\vp}_t(v^{(\ell)})
\right|^2
\right)
\mathop{\longrightarrow}\limits_{\vp \to 0} 0.
$$
\end{lemma}

\begin{proof}
Let $(v^{(\ell)})_{\ell \ge 1}$ be a sequence of real numbers. The coupling of the two processes is defined using the same perimeter process $P$, the same exponential random variables for $\widetilde{P}$ but the random variables $\overline{U}^{(\ell),T}$ defining $\overline{X}^{(\ell),T,\vp}_t(v^{(\ell)})$ are defined inductively by the relations
\begin{align}
\overline{U}_{T-t}^{(\ell),T,\vp}&= \left\{ \overline{X}^{(\ell),T,\vp}_{t-}(v^{(\ell)}) - X^{(\ell),T}_{t-}(v^{(\ell)}) + U^{(\ell)}_{T-t}\right\}, \\
\overline{N}^{(\ell),T,\vp}_{\vert [0,t]\times (\R^*_+\setminus \{1\})\times [0,1]} 
&= \sum_{s\in [T-t,T]; \ \widetilde{P}(s)/\widetilde{P}(s-) \in \R_+^*\setminus (1-\vp, 1+\vp)} \delta_{(T-s,\widetilde{P}(s)/\widetilde{P}(s-), \overline{U}^{(\ell),T,\vp}_{s})}, \\
\overline{X}^{(\ell),T,\vp}_t(v^{(\ell)}) &= v+ \int_0^t \int_0^\infty \int_0^1 g^{(\ell),T}_s(\overline{X}^{(\ell),T,\vp}_{s-}(v^{(\ell)}),z,u) \overline{N}^{(\ell), T}_\vp(ds,dz,du),
\end{align}
for all $t\in [0,T]$ such that $\Delta \widetilde{P}(t)\neq 0$. Then one can see that $\overline{X}^{(\ell),T,\vp}(v^{(\ell)})$ has the same law as ${X}^{(\ell),T,\vp}(v^{(\ell)})$ and that by construction, for all $t\in [0,T]$ such that $\Delta \widetilde{P}(t)\neq 0$, for all $z \in \R_+^*\setminus \left\{1\right\}$, 
\begin{equation}\label{eq egalite couplage}
g^{(\ell),T}_t(\overline{X}^{(\ell),T,\vp}_{t-}(v^{(\ell)}), z, \overline{U}_{T-t}^{(\ell),T,\vp}) = g^{(\ell),T}_t(X^{(\ell),T}_{t-}(v^{(\ell)}), z, U^{(\ell)}_{T-t}).
\end{equation}
Using the coupling, we see that
\begin{align*}
\E^{(\ell)} 
&\left(\sup_{t\in [0,T]} 
\left|
X^{(\ell),T}_t(v^{(\ell)})
-
\overline{X}^{(\ell),T,\vp}_t(v^{(\ell)})
\right|^2
\right)
\\
=
&\E^{(\ell)} 
\left(\sup_{t\in [0,T]} 
\left|
\int_0^t \int_{1-\vp}^{1+\vp} \int_0^1
g^{(\ell),T}_s(X^{(\ell),T}_{s-}(v^{(\ell)}),z,u)
N^{(\ell),T}(ds, dz, du)
\right|^2
\right)
\qquad
\text{by } \eqref{eq egalite couplage}.
\end{align*}
Moreover, thanks to Lemma \ref{lemme martingale}, the process
$$
(M_t^{(\ell),T,\vp})_{t \in [0,T]}
=
\left(
\int_{0}^t
\int_{1-\vp}^{1+\vp}
\int_0^1
g^{(\ell),T}_s(X^{(\ell),T}_{s-}(v^{(\ell)}),z,u) {N}^{(\ell),T} (ds, dz, du)
\right)_{t \in [0,T]}
$$
is a martingale. So by Doob's martingale inequality,
$$
\E^{(\ell)}
\left(
\sup_{t \in [0,T]}
(M^{(\ell),T,\vp}_t)^2
\right)
\le  4 \sup_{t \in [0,T]} \E^{(\ell)} \left( (M^{(\ell),T,\vp}_t)^2\right).
$$
Moreover, by \eqref{eq majoration g discret}, one can {bound from above} for all $t \in [0,T]$,
\begin{align*}
\E^{(\ell)} \left( (M^{(\ell),T,\vp}_t)^2 \right)
&=\E^{(\ell)} \left(\int_{0}^t \int_{1-\vp}^{1+\vp} \int_0^1
g^{(\ell),T}_s(X^{(\ell),T}_{s-}(v^{(\ell)}),z,u)^2 
{N}^{(\ell),T}(ds,dz,du)
\right) \enskip\text{(quadratic variation)}\\
&\le 9\E^{(\ell)} \left(\int_{0}^t \int_{1-\vp}^{1+\vp} \int_0^1
\vert z-1\vert^2 
{N}^{(\ell),T}(ds,dz,du)
\right) \\
&\le 9 \E^{(\ell)}
\left(
\sum_{s \in[0,T] } {\bf 1}_{ \widetilde{P}(s)/\widetilde{P}(s-) \in [1-\vp,1+\vp]\setminus \{1\} }
\left(\frac{\Delta \widetilde{P}(s)}{\widetilde{P}(s-)}\right)^2 
\right).
\end{align*}
Using Lemma \ref{lemme petits sauts de P} one concludes that
\begin{equation}\label{eq majoration unif de la martingale}
	\sup_{\ell \ge 1} \sup_{t \in [0,T]} \E^{(\ell)} \left( (M^{(\ell),T,\vp}_t)^2\right)
	\mathop{\longrightarrow}\limits_{\vp \to 0} 0.
\end{equation}
This ends the proof.
\end{proof}

\end{document}
